\chapter{Photosynthesis}\label{ch:g10:photosynthesis}

Sunlight is the energy income of almost every ecosystem. \omterm{def:g10:photosynthesis:def}{Photosynthesis}
captures that energy in chemical bonds of sugars, releasing oxygen as a
by-product. This chapter places the process in the \omterm{def:g10:photosynthesis:chloroplast}{chloroplast}, states
the overall equation, and separates the \omterm{def:g10:photosynthesis:light}{light-dependent reactions} from
the \omterm{def:g10:photosynthesis:calvin}{Calvin cycle} --- enough structure to see why light, $\mathrm{CO_2}$
and temperature all matter.

\section{Where and why}

\begin{definition}[Photosynthesis]\label{def:g10:photosynthesis:def}
\emph{Photosynthesis}\index{photosynthesis} is the process by which
photoautotrophs (plants, algae and cyanobacteria) convert light energy
into chemical energy stored in organic molecules, using carbon dioxide
and water as raw materials and releasing oxygen (in oxygenic
photosynthesis).
\end{definition}

\begin{proposition}[Overall equation]\label{prop:g10:photosynthesis:equation}
Oxygenic \omterm{def:g10:photosynthesis:def}{photosynthesis} is summarised by
\[
6\,\mathrm{CO_2} + 6\,\mathrm{H_2O}
\xrightarrow{\text{light}}
\mathrm{C_6H_{12}O_6} + 6\,\mathrm{O_2}
\]
(or, more generally, $\mathrm{CO_2} + \mathrm{H_2O} \to
(\mathrm{CH_2O}) + \mathrm{O_2}$). Glucose stands for \omterm{def:g10:molecules-of-life:carbs}{carbohydrate}
product; real \omterm{def:g10:the-cell:cell}{cells} make intermediates that feed sucrose, starch and
other pathways.
\end{proposition}
\admitted

\begin{remark}[Who depends on it]\label{rem:g10:photosynthesis:who}
Consumers and most decomposers live on energy originally fixed by
\omterm{def:g10:photosynthesis:def}{photosynthesis} (or, rarely, chemosynthesis). Atmospheric $\mathrm{O_2}$
is largely a photosynthetic legacy. The process is therefore both a \omterm{def:g10:the-cell:cell}{cell}
pathway and a planetary service.
\end{remark}

\section{The chloroplast}

\begin{definition}[Chloroplast structure]\label{def:g10:photosynthesis:chloroplast}
The \emph{chloroplast}\index{chloroplast} is the \omterm{def:g10:the-cell:prok-euk}{organelle} of
\omterm{def:g10:photosynthesis:def}{photosynthesis} in \omterm{def:g10:the-cell:prok-euk}{eukaryotes}. It has a double outer membrane enclosing
the \emph{stroma}\index{stroma} (fluid). Stacked membrane sacs called
\emph{thylakoids}\index{thylakoid} form \emph{grana}\index{granum};
their membranes hold pigments and the light-reaction machinery. The
stroma holds \omterm{def:g10:membranes-enzymes:enzyme}{enzymes} of the \omterm{def:g10:photosynthesis:calvin}{Calvin cycle}.
\end{definition}

\begin{definition}[Chlorophyll]\label{def:g10:photosynthesis:chlorophyll}
\emph{Chlorophyll}\index{chlorophyll} is the main photosynthetic pigment:
it absorbs mainly violet--blue and red light and reflects green ---
hence green leaves. Accessory pigments (carotenoids and others) broaden
the harvest and protect against excess light.
\end{definition}

\begin{omfigure}
\begin{tikzpicture}[font=\footnotesize]
  \draw[thick, omMeth, fill=omMeth!10] (0,0) ellipse (3.4 and 1.9);
  \draw[thick, omMeth] (0,0) ellipse (3.15 and 1.65);
  \foreach \y in {-0.9,-0.15,0.6} {
    \draw[thick, omDef, fill=omDef!25] (-1.6,\y) ellipse (1.1 and 0.28);
    \draw[thick, omDef, fill=omDef!25] (-1.6,\y+0.22) ellipse (1.1 and 0.28);
  }
  \node[omDef] at (-1.6,1.25) {grana (thylakoids)};
  \node[omProp] at (1.5,0.2) {stroma};
  \node at (0,-2.25) {chloroplast (schematic)};
\end{tikzpicture}

{\small Light reactions on \omterm{def:g10:photosynthesis:chloroplast}{thylakoid} membranes; sugar-building \omterm{def:g10:photosynthesis:calvin}{Calvin
cycle} in the \omterm{def:g10:photosynthesis:chloroplast}{stroma}.}
\end{omfigure}

\section{Two stages}

\begin{definition}[Light-dependent reactions]\label{def:g10:photosynthesis:light}
In the \emph{light-dependent reactions}\index{light-dependent reactions}
(\omterm{def:g10:photosynthesis:chloroplast}{thylakoid} membrane), light energy excites electrons in \omterm{def:g10:photosynthesis:chlorophyll}{chlorophyll}.
Water is split (\emph{photolysis}\index{photolysis}), releasing
$\mathrm{O_2}$, protons and electrons. Electron transport produces ATP
and the reduced carrier NADPH.
\end{definition}

\begin{definition}[Calvin cycle]\label{def:g10:photosynthesis:calvin}
The \emph{Calvin cycle}\index{Calvin cycle} (\omterm{def:g10:photosynthesis:chloroplast}{stroma}) uses ATP and NADPH
to fix $\mathrm{CO_2}$ into \omterm{def:g10:molecules-of-life:carbs}{carbohydrate}. Ribulose bisphosphate (RuBP)
accepts $\mathrm{CO_2}$ (catalysed by Rubisco); products are rearranged
to form sugars and regenerate RuBP. Light is not required
\emph{directly}, but the cycle stops when ATP/NADPH run out in the dark.
\end{definition}

\begin{omfigure}
\begin{tikzpicture}[font=\small,
  box/.style={draw, rounded corners, minimum width=2.6cm,
    minimum height=1.0cm, align=center}]
  \node[box, fill=omDef!15] (light) at (0,0) {light\\reactions};
  \node[box, fill=omMeth!15] (cal) at (5.2,0) {Calvin\\cycle};
  \draw[-Stealth, thick, omThm] (-2.2,0.8) -- (light.west)
    node[midway, above, font=\footnotesize] {light};
  \draw[-Stealth, thick] (-2.2,-0.3) -- (light.west)
    node[midway, below, font=\footnotesize] {$\mathrm{H_2O}$};
  \draw[-Stealth, thick, omProp] (light.east) -- (cal.west)
    node[midway, above, font=\footnotesize] {ATP, NADPH};
  \draw[-Stealth, thick] (light.south) -- ++(0,-0.7)
    node[below, font=\footnotesize] {$\mathrm{O_2}$};
  \draw[-Stealth, thick] (5.2,1.4) -- (cal.north)
    node[midway, right, font=\footnotesize] {$\mathrm{CO_2}$};
  \draw[-Stealth, thick, omMeth] (cal.east) -- ++(1.3,0)
    node[right, font=\footnotesize] {sugars};
\end{tikzpicture}

{\small Energy and electrons flow from light reactions to the \omterm{def:g10:photosynthesis:calvin}{Calvin
cycle}; $\mathrm{O_2}$ exits from water splitting.}
\end{omfigure}

\begin{proposition}[Oxygen comes from water]\label{prop:g10:photosynthesis:oxygen-source}
The $\mathrm{O_2}$ released by oxygenic \omterm{def:g10:photosynthesis:def}{photosynthesis} originates from
water, not from $\mathrm{CO_2}$ --- shown historically by isotope
labelling \omterm{def:g10:scientific-biology:experiment}{experiments}.
\end{proposition}
\admitted

\begin{example}[Leaf as factory]\label{ex:g10:photosynthesis:leaf}
$\mathrm{CO_2}$ enters through stomata; water arrives via xylem; light
hits mesophyll \omterm{def:g10:photosynthesis:chloroplast}{chloroplasts}; sugars are exported in phloem or stored as
starch. Closing stomata to save water can starve the \omterm{def:g10:photosynthesis:calvin}{Calvin cycle} of
$\mathrm{CO_2}$.
\end{example}

\begin{omfigure}
\begin{tikzpicture}[font=\footnotesize]
  % leaf cross-section schematic
  \draw[thick] (0,0) rectangle (7.2,2.8);
  \fill[omDef!15] (0.15,2.2) rectangle (7.05,2.65);
  \node at (3.6,2.42) {upper epidermis};
  \fill[omMeth!20] (0.15,1.35) rectangle (7.05,2.2);
  \foreach \x in {0.7,1.8,2.9,4.0,5.1,6.2}
    \fill[omDef!50] (\x,1.55) ellipse (0.28 and 0.28);
  \node at (3.6,1.9) {palisade mesophyll (many chloroplasts)};
  \fill[omProp!15] (0.15,0.55) rectangle (7.05,1.35);
  \foreach \x/\y in {1.0/0.95,2.2/0.8,3.4/1.0,4.6/0.75,5.8/0.95}
    \fill[omDef!40] (\x,\y) circle [radius=0.18cm];
  \node at (3.6,1.15) {spongy mesophyll + air spaces};
  \fill[omThm!15] (0.15,0.15) rectangle (7.05,0.55);
  \node at (3.6,0.35) {lower epidermis};
  \draw[thick, omExo] (5.6,0.15) -- (5.6,-0.05) -- (6.2,-0.05) -- (6.2,0.15);
  \node at (5.9,-0.35) {stoma};
\end{tikzpicture}

{\small Leaf cross-section (schematic): gas exchange at stomata; most
\omterm{def:g10:photosynthesis:chloroplast}{chloroplasts} in mesophyll, especially the packed palisade layer.}
\end{omfigure}

\begin{proposition}[Absorption spectrum idea]
\label{prop:g10:photosynthesis:absorption}
An \emph{absorption spectrum}\index{absorption spectrum} plots how strongly a
pigment absorbs different wavelengths. \omterm{def:g10:photosynthesis:chlorophyll}{Chlorophylls} absorb mainly blue and red
light and absorb green poorly --- hence green reflection. The
\emph{action spectrum}\index{action spectrum} of \omterm{def:g10:photosynthesis:def}{photosynthesis} (rate versus
wavelength) roughly follows pigment absorption, with accessory pigments
broadening useful light.
\end{proposition}
\admitted

\begin{omfigure}
\begin{tikzpicture}[font=\small]
  \draw[thick, -Stealth] (0,0) -- (7.2,0) node[right] {wavelength};
  \draw[thick, -Stealth] (0,0) -- (0,3.2) node[above] {absorption};
  \draw[omDef, very thick, smooth]
    plot coordinates {(0.5,0.3) (1.2,2.6) (2.0,1.8) (2.8,0.5)
      (3.8,0.35) (5.0,0.4) (5.8,2.3) (6.5,0.4)};
  \node[font=\footnotesize] at (1.4,2.9) {blue};
  \node[font=\footnotesize] at (3.5,0.7) {green};
  \node[font=\footnotesize] at (5.8,2.7) {red};
  \node[align=center, font=\footnotesize] at (3.5,-0.7)
    {chlorophyll absorption sketch (not to scale)};
\end{tikzpicture}

{\small \omterm{def:g10:photosynthesis:chlorophyll}{Chlorophyll} absorbs blue and red strongly and green weakly ---
matching why leaves look green under white light.}
\end{omfigure}

\begin{remark}[C$_3$, C$_4$ and CAM --- one-paragraph comparison]
\label{rem:g10:photosynthesis:C3C4CAM}
Most temperate plants are \emph{C$_3$}\index{C3 plant}: Rubisco fixes
$\mathrm{CO_2}$ directly in mesophyll, which is efficient in cool moist
climates but vulnerable to photorespiration when stomata close in heat.
\emph{C$_4$}\index{C4 plant} plants (many tropical grasses) first fix carbon
into a four-carbon compound in mesophyll and release $\mathrm{CO_2}$ near
Rubisco in bundle-sheath \omterm{def:g10:the-cell:cell}{cells}, concentrating $\mathrm{CO_2}$ and reducing
photorespiration at a metabolic cost. \emph{CAM}\index{CAM plant} plants
(e.g.\ many succulents) open stomata at night to store $\mathrm{CO_2}$ as acid
and close them by day, trading \omterm{def:g10:scientific-biology:properties}{growth} rate for extreme water conservation.
All three still use the \omterm{def:g10:photosynthesis:calvin}{Calvin cycle}; they differ in \emph{when} and
\emph{where} $\mathrm{CO_2}$ is first captured.
\end{remark}

\section{Factors affecting the rate}

\begin{definition}[Limiting factors]\label{def:g10:photosynthesis:limiting}
A \emph{limiting factor}\index{limiting factor} is an environmental
variable that, when increased, raises the rate because it is in
shortest supply relative to demand. For \omterm{def:g10:photosynthesis:def}{photosynthesis}, common candidates
are light intensity, $\mathrm{CO_2}$ concentration and temperature
(\omterm{def:g10:membranes-enzymes:enzyme}{enzyme} kinetics, including Rubisco).
\end{definition}

\begin{method}[Interpreting rate curves]\label{met:g10:photosynthesis:curves}
\begin{enumerate}
  \item At low light, rate rises roughly with intensity; light limits.
  \item A plateau means another factor (or the machinery's maximum)
        limits.
  \item Raising $\mathrm{CO_2}$ or temperature may lift the plateau if
        that factor was limiting --- until a new limit appears.
  \item Very high temperature can denature \omterm{def:g10:membranes-enzymes:enzyme}{enzymes} and collapse the rate.
\end{enumerate}
\end{method}

\begin{omfigure}
\begin{tikzpicture}[font=\small]
  \begin{axis}[
    width=8cm, height=5cm,
    axis lines=left,
    xlabel={light intensity},
    ylabel={photosynthesis rate},
    xmin=0, xmax=10, ymin=0, ymax=5,
    xtick=\empty, ytick=\empty,
    legend style={font=\footnotesize, at={(0.97,0.35)}, anchor=east},
  ]
    \addplot[thick, omDef, domain=0:10, samples=60] {3.2*x/(x+1.5)};
    \addplot[thick, omThm, domain=0:10, samples=60] {4.3*x/(x+1.5)};
    \legend{low $\mathrm{CO_2}$, high $\mathrm{CO_2}$}
  \end{axis}
\end{tikzpicture}

{\small Light curves: both rise then plateau; higher $\mathrm{CO_2}$
raises the plateau when carbon was limiting.}
\end{omfigure}

\begin{proposition}[Compensation point (idea)]\label{prop:g10:photosynthesis:compensation}
At the light \emph{compensation point}\index{compensation point},
photosynthetic $\mathrm{O_2}$ production (or $\mathrm{CO_2}$ uptake)
balances respiratory consumption. Below it, the plant is a net consumer
of \omterm{def:g10:molecules-of-life:carbs}{carbohydrate}; above it, a net producer.
\end{proposition}
\admitted

\section{Exercises}

\begin{exercise}[$\star$]\label{exo:g10:photosynthesis:1}
Write the overall word equation and the balanced symbolic equation for
oxygenic \omterm{def:g10:photosynthesis:def}{photosynthesis}.
\end{exercise}

\begin{exercise}[$\star$]\label{exo:g10:photosynthesis:2}
Where in the \omterm{def:g10:photosynthesis:chloroplast}{chloroplast} do (a) the \omterm{def:g10:photosynthesis:light}{light-dependent reactions} and
(b) the \omterm{def:g10:photosynthesis:calvin}{Calvin cycle} occur?
\end{exercise}

\begin{exercise}[$\star$]\label{exo:g10:photosynthesis:3}
Why do most leaves look green under white light?
\end{exercise}

\begin{exercise}[$\star$]\label{exo:g10:photosynthesis:4}
Name the products of the \omterm{def:g10:photosynthesis:light}{light-dependent reactions} that the \omterm{def:g10:photosynthesis:calvin}{Calvin cycle}
consumes. What gas is released when water is split?
\end{exercise}

\begin{exercise}[$\star$]\label{exo:g10:photosynthesis:5}
List three environmental factors that can limit the rate of
\omterm{def:g10:photosynthesis:def}{photosynthesis} and state one way each might be raised in a greenhouse.
\end{exercise}

\begin{exercise}[$\star\star$]\label{exo:g10:photosynthesis:6}
Explain why the \omterm{def:g10:photosynthesis:calvin}{Calvin cycle} is sometimes called the ``light-independent''
reactions, and why that name can mislead students.
\end{exercise}

\begin{exercise}[$\star\star$]\label{exo:g10:photosynthesis:7}
A graph of \omterm{def:g10:photosynthesis:def}{photosynthesis} rate versus light intensity plateaus. What does
the plateau mean? How could you test whether $\mathrm{CO_2}$ is the
\omterm{def:g10:photosynthesis:limiting}{limiting factor} at that light level?
\end{exercise}

\begin{exercise}[$\star\star$]\label{exo:g10:photosynthesis:8}
Isotope \omterm{def:g10:scientific-biology:experiment}{experiment} idea: if water is labelled with $^{18}\mathrm{O}$ but
$\mathrm{CO_2}$ is not, where does the label appear in the products?
What if only $\mathrm{CO_2}$ is labelled?
\end{exercise}

\begin{exercise}[$\star\star$]\label{exo:g10:photosynthesis:9}
Stomata close on a hot dry day. Predict effects on $\mathrm{CO_2}$ entry,
transpiration and short-term sugar production. Trade-off?
\end{exercise}

\begin{exercise}[$\star\star\star$]\label{exo:g10:photosynthesis:10}
Design a floating-disk or Elodea bubble \omterm{def:g10:scientific-biology:experiment}{experiment} to compare
\omterm{def:g10:photosynthesis:def}{photosynthesis} rate at two light intensities. Name independent,
dependent and three controlled variables.
\end{exercise}

\begin{exercise}[$\star\star\star$]\label{exo:g10:photosynthesis:11}
A forest understory plant has large, thin leaves and a low light
\omterm{prop:g10:photosynthesis:compensation}{compensation point}. Explain how these traits suit a shady habitat.
\end{exercise}

\begin{exercise}[$\star\star$]\label{exo:g10:photosynthesis:12}
Sketch or describe an \omterm{prop:g10:photosynthesis:absorption}{absorption spectrum} for \omterm{def:g10:photosynthesis:chlorophyll}{chlorophyll} and link it to
leaf colour. Why does an \omterm{prop:g10:photosynthesis:absorption}{action spectrum} of \omterm{def:g10:photosynthesis:def}{photosynthesis} not look
identical to a pure \omterm{def:g10:photosynthesis:chlorophyll}{chlorophyll} \omterm{prop:g10:photosynthesis:absorption}{absorption spectrum}?
\end{exercise}

\begin{exercise}[$\star\star\star$]\label{exo:g10:photosynthesis:13}
In one short paragraph each, contrast C$_3$, C$_4$ and \omterm{rem:g10:photosynthesis:C3C4CAM}{CAM} strategies:
when/where $\mathrm{CO_2}$ is first fixed, and the climate problem each
mainly addresses. What do all three share?
\end{exercise}

\section{Problem: Powering a biosphere bottle}

\begin{problem}[{Weekend problem --- balance light, gas and biomass in
a sealed micro-ecosystem: predict day--night gas swings and diagnose
failure modes}]
\label{pb:g10:photosynthesis:1}
You build a sealed transparent bottle with aquatic plants, pond water
and a few snails, then leave it under a day/night light cycle.

\medskip
\noindent\textbf{Part I --- Daytime budget.}
\begin{enumerate}
  \item Write the net photosynthetic equation and state what happens to
        $\mathrm{O_2}$ and $\mathrm{CO_2}$ levels by afternoon if light
        is bright.
  \item Where in plant \omterm{def:g10:the-cell:cell}{cells} is $\mathrm{O_2}$ produced?
\end{enumerate}

\medskip
\noindent\textbf{Part II --- Night.}
\begin{enumerate}
  \item Plants and snails respire. Predict overnight trends in
        $\mathrm{O_2}$ and $\mathrm{CO_2}$.
  \item Why might the snails die first if the bottle is left in continuous
        darkness?
\end{enumerate}

\medskip
\noindent\textbf{Part III --- Troubleshooting.}
\begin{enumerate}
  \item Algae coat the glass and light inside falls. Predict the effect
        on net \omterm{def:g10:photosynthesis:def}{photosynthesis} of the rooted plants.
  \item You open the bottle daily ``for air''. Explain how that changes
        the \omterm{def:g10:scientific-biology:experiment}{experiment} as a model of a closed system.
  \item Propose two measurements (qualitative or quantitative) that would
        show the plants are photosynthesising.
\end{enumerate}
\end{problem}
